\section{Derived limits and the exact truncation-recovery obstruction}\label{proofs-10435-10686-derived-limits-and-the-exact-truncation-recovery-obstruction}

Private source-order review of original derived.tex 10435–\allowbreak{}10686. Authority:\allowbreak{} Stacks a04446e5\allowbreak{}7ec1fbc2\allowbreak{}52a871af\allowbreak{}cec7752f\allowbreak{}b2807b14. Receiving public edition:\allowbreak{} 95a51593\allowbreak{}aceb247a\allowbreak{}c7a18311\allowbreak{}1678c3a0\allowbreak{}f8b8f4b5. The operation and source-reading records bind the files. The two consequences in Sections 6–\allowbreak{}7 are separate editorial results; the original hypotheses and theorem statements are retained. No novelty is claimed.

Human source:\allowbreak{} \href{https://github.com/stacks/stacks-project/blob/a04446e57ec1fbc252a871afcec7752fb2807b14/derived.tex\#L10435}{The Stacks Project,\allowbreak{} original Derived Categories TeX}. The receiving source is \href{https://github.com/stacks/stacks-project/blob/a04446e57ec1fbc252a871afcec7752fb2807b14/more-algebra.tex}{More on Algebra,\allowbreak{} original TeX},\allowbreak{} lemmas named lemma-pro-isomorphism-iso-Rlim,\allowbreak{} lemma-map-into-Rlim,\allowbreak{} and lemma-pro-isomorphism-Rlim.

\subsection{1. Products and the defining triangle}\label{proofs-10435-10686-products-and-the-defining-triangle}

For a tower with transition maps f\_\allowbreak{}\{n\allowbreak{}+1\}:\allowbreak{}K\_\allowbreak{}\{n\allowbreak{}+1\}\allowbreak{}→K\_\allowbreak{}n,\allowbreak{} let P\allowbreak{}=∏\_\allowbreak{}n K\_\allowbreak{}n with projections π\_\allowbreak{}n and define Δ:\allowbreak{}P\allowbreak{}→P by

π\_\allowbreak{}n Δ\allowbreak{}=π\_\allowbreak{}n-f\_\allowbreak{}\{n\allowbreak{}+1\}π\_\allowbreak{}\{n\allowbreak{}+1\}.

This is a map specified by its product components,\allowbreak{} not an assumption that the K\_\allowbreak{}n have elements. A homotopy limit is a distinguished triangle

C --j-\allowbreak{}-> P --Δ-\allowbreak{}-> P --w-\allowbreak{}-> C[1]\allowbreak{}. (1.1)\allowbreak{}

Existence follows by completing Δ to a triangle and rotating. Two such triangles have an isomorphism respecting j,\allowbreak{} Δ and w,\allowbreak{} by TR3 and the two-isomorphisms criterion for maps of triangles. This need not choose a unique map between their first objects.

Suppose A is abelian with exact countable products. The termwise product of complexes has H\textasciicircum{}p(∏K\_\allowbreak{}n)\allowbreak{}\allowbreak{}=∏H\textasciicircum{}p(K\_\allowbreak{}n)\allowbreak{}:\allowbreak{} apply exact products to the short exact sequences for kernels,\allowbreak{} images and cohomology. Thus products preserve quasi-isomorphisms. For maps L\allowbreak{}→K\_\allowbreak{}n in D(A)\allowbreak{},\allowbreak{} choose right roofs L\allowbreak{}→f\_\allowbreak{}n M\_\allowbreak{}n←s\_\allowbreak{}n K\_\allowbreak{}n,\allowbreak{} where s\_\allowbreak{}n is a quasi-isomorphism. Then s\allowbreak{}=∏s\_\allowbreak{}n is a quasi-isomorphism,\allowbreak{} and the maps f\_\allowbreak{}n give f:\allowbreak{}L\allowbreak{}→∏M\_\allowbreak{}n. The right roof s\textasciicircum{}(-1)\allowbreak{}f gives the desired product map.

Here is the uniqueness omitted in the source. Put P\allowbreak{}=∏K\_\allowbreak{}n termwise,\allowbreak{} and suppose γ:\allowbreak{}L\allowbreak{}→P has zero composite with every projection. Write γ\allowbreak{}=s\textasciicircum{}(-1)\allowbreak{}f with s:\allowbreak{}P\allowbreak{}→M a quasi-isomorphism. For each projection p\_\allowbreak{}n:\allowbreak{}P\allowbreak{}→K\_\allowbreak{}n,\allowbreak{} right fraction calculus gives u\_\allowbreak{}n:\allowbreak{}M\allowbreak{}→N\_\allowbreak{}n and a quasi-isomorphism t\_\allowbreak{}n:\allowbreak{}K\_\allowbreak{}n\allowbreak{}→N\_\allowbreak{}n with u\_\allowbreak{}n s\allowbreak{}=t\_\allowbreak{}n p\_\allowbreak{}n in K(A)\allowbreak{}. Since u\_\allowbreak{}n f is zero in D(A)\allowbreak{},\allowbreak{} the zero criterion for localization allows postcomposition by a further quasi-isomorphism N\_\allowbreak{}n\allowbreak{}→N'\_\allowbreak{}n to make u\_\allowbreak{}n f zero in K(A)\allowbreak{}. Use these refined u\_\allowbreak{}n and t\_\allowbreak{}n.

Form U:\allowbreak{}M\allowbreak{}→∏N'\_\allowbreak{}n and T\allowbreak{}=∏t\_\allowbreak{}n:\allowbreak{}P\allowbreak{}→∏N'\_\allowbreak{}n. Coordinate homotopies combine to U s\allowbreak{}=T and U f\allowbreak{}=0 in K(A)\allowbreak{},\allowbreak{} by the degreewise product property. T and s are quasi-isomorphisms,\allowbreak{} so U is one. It follows that γ\allowbreak{}=(U s)\allowbreak{}\textasciicircum{}(-1)\allowbreak{}U f\allowbreak{}=0. This proves the product property in D(A)\allowbreak{} without assuming it in the uniqueness argument.

With enough injectives and merely existing countable products,\allowbreak{} a tower whose individual objects lie in D\textasciicircum{}\allowbreak{}+(A)\allowbreak{} also has a derived product:\allowbreak{} choose a bounded below injective representative I\_\allowbreak{}n for each object. Every I\_\allowbreak{}n is K-injective,\allowbreak{} and the full Hom-complex product argument in PART14 proves ∏I\_\allowbreak{}n is K-injective and represents the product in D(A)\allowbreak{}. No common lower bound is required. Its difference map gives (1.1)\allowbreak{}. This proves the source\textquotesingle s second existence criterion without adding exactness of products.

\subsection{2. Maps into a homotopy limit}\label{proofs-10435-10686-maps-into-a-homotopy-limit}

Use the shifted boundary k\allowbreak{}=-w[-1]\allowbreak{}:\allowbreak{}P[-1]\allowbreak{}\allowbreak{}→C. With the source rotation convention this is the first arrow of the distinguished triangle P[-1]\allowbreak{}\allowbreak{}→C\allowbreak{}→P\allowbreak{}→P. Applying Hom(T,\allowbreak{}-)\allowbreak{} gives the exact segment

Hom(T,\allowbreak{}P[-1]\allowbreak{})\allowbreak{} --Δ[-1]\allowbreak{}\_\allowbreak{}*-\allowbreak{}-> Hom(T,\allowbreak{}P[-1]\allowbreak{})\allowbreak{} --k\_\allowbreak{}*-\allowbreak{}-> Hom(T,\allowbreak{}C)\allowbreak{} --j\_\allowbreak{}*-\allowbreak{}-> Hom(T,\allowbreak{}P)\allowbreak{} --Δ\_\allowbreak{}*-\allowbreak{}-> Hom(T,\allowbreak{}P)\allowbreak{}.

Representable Hom preserves products. Thus Δ\_\allowbreak{}* is precisely the difference map on ∏Hom(T,\allowbreak{}K\_\allowbreak{}n)\allowbreak{},\allowbreak{} and the preceding map is the difference map on ∏Hom(T,\allowbreak{}K\_\allowbreak{}n[-1]\allowbreak{})\allowbreak{}. Consequently

0\allowbreak{}→coker Δ\_\allowbreak{}\{Hom(T,\allowbreak{}K\_\allowbreak{}n[-1]\allowbreak{})\allowbreak{}\}\allowbreak{}→Hom(T,\allowbreak{}C)\allowbreak{} \allowbreak{}→lim\_\allowbreak{}n Hom(T,\allowbreak{}K\_\allowbreak{}n)\allowbreak{}\allowbreak{}→0. (2.1)\allowbreak{}

The injection sends a family b:\allowbreak{}T\allowbreak{}→P[-1]\allowbreak{} to k b\allowbreak{}=-w[-1]\allowbreak{}b,\allowbreak{} with its indicated sign. The right map is given by all π\_\allowbreak{}n j. The cokernel is the R\^{}1 lim of this abelian-group tower,\allowbreak{} by the explicit two-term complex from More on Algebra. In particular compatible maps T\allowbreak{}→K\_\allowbreak{}n lift to T\allowbreak{}→C,\allowbreak{} and their possible lifts form a torsor under the left group. Nothing here requires exact products in the underlying abelian category when the needed product already exists in D.

\subsection{\texorpdfstring{3. Pro-morphisms:\allowbreak{} exact finite sums and independence}{3.  exact finite sums and independence}}\label{proofs-10435-10686-pro-morphisms-exact-finite-sums-and-independence}

Write t\_\allowbreak{}\{r,\allowbreak{}s\}:\allowbreak{}K\_\allowbreak{}s\allowbreak{}→K\_\allowbreak{}r for the original transition composite,\allowbreak{} with t\_\allowbreak{}\{r,\allowbreak{}r\}\allowbreak{}=id. Let g denote transitions in a second tower L. A pro-morphism may be represented by strictly increasing indices m\_\allowbreak{}n and maps a\_\allowbreak{}n:\allowbreak{}K\_\allowbreak{}\{m\_\allowbreak{}n\}\allowbreak{}→L\_\allowbreak{}n satisfying g\_\allowbreak{}\{n\allowbreak{}+1\}a\_\allowbreak{}\{n\allowbreak{}+1\}\allowbreak{}=a\_\allowbreak{}n t\_\allowbreak{}\{m\_\allowbreak{}n,\allowbreak{}m\_\allowbreak{}\{n\allowbreak{}+1\}\}. To obtain such a representation,\allowbreak{} choose successive representatives of the inverse-limit-of-colimits definition of a pro-morphism. At stage n\allowbreak{}+1,\allowbreak{} choose an index larger than m\_\allowbreak{}n and large enough that the required equality with the n-th representative holds. That equality is available after a sufficiently late transition by the definition of equality in the filtered colimit. Earlier equalities then follow by composition. This also permits m\_\allowbreak{}n≥n.

The source\textquotesingle s two product maps have exact components

π\_\allowbreak{}n a'\allowbreak{}=a\_\allowbreak{}n π\_\allowbreak{}\{m\_\allowbreak{}n\},\allowbreak{} π\_\allowbreak{}n a''\allowbreak{}=Σ\_\allowbreak{}\{m\_\allowbreak{}n≤r<m\_\allowbreak{}\{n\allowbreak{}+1\}\}a\_\allowbreak{}n t\_\allowbreak{}\{m\_\allowbreak{}n,\allowbreak{}r\}π\_\allowbreak{}r. (3.1)\allowbreak{}

The sum in each component is finite. Telescoping proves

π\_\allowbreak{}n a''Δ\_\allowbreak{}K \allowbreak{}=a\_\allowbreak{}n π\_\allowbreak{}\{m\_\allowbreak{}n\}-a\_\allowbreak{}n t\_\allowbreak{}\{m\_\allowbreak{}n,\allowbreak{}m\_\allowbreak{}\{n\allowbreak{}+1\}\}π\_\allowbreak{}\{m\_\allowbreak{}\{n\allowbreak{}+1\}\} \allowbreak{}=π\_\allowbreak{}n Δ\_\allowbreak{}L a\textquotesingle.

Thus Δ\_\allowbreak{}L a'\allowbreak{}=a''Δ\_\allowbreak{}K. TR3,\allowbreak{} after rotation,\allowbreak{} gives a map C\_\allowbreak{}K\allowbreak{}→C\_\allowbreak{}L and a map of triangles with product components a\textquotesingle{} and a\textquotesingle\textquotesingle. This verifies every transition factor in original 10539. The spelling correction in unit 0886 changes none of them.

The construction induces well-defined maps on kernel and cokernel of the abelian-group difference maps. Here are the full comparison homotopies,\allowbreak{} including the omitted representation-independence step. If m\_\allowbreak{}n is enlarged to M\_\allowbreak{}n≥m\_\allowbreak{}n and b\_\allowbreak{}n\allowbreak{}=a\_\allowbreak{}n t\_\allowbreak{}\{m\_\allowbreak{}n,\allowbreak{}M\_\allowbreak{}n\},\allowbreak{} define a map H:\allowbreak{}∏K\_\allowbreak{}r\allowbreak{}→∏L\_\allowbreak{}n by

π\_\allowbreak{}n H\allowbreak{}=Σ\_\allowbreak{}\{m\_\allowbreak{}n≤r<M\_\allowbreak{}n\}a\_\allowbreak{}n t\_\allowbreak{}\{m\_\allowbreak{}n,\allowbreak{}r\}π\_\allowbreak{}r.

The same finite telescoping gives

a'-b'\allowbreak{}=HΔ\_\allowbreak{}K,\allowbreak{} a''-b''\allowbreak{}=Δ\_\allowbreak{}L H. (3.2)\allowbreak{}

For the second equality,\allowbreak{} expand π\_\allowbreak{}n H-g\_\allowbreak{}\{n\allowbreak{}+1\}π\_\allowbreak{}\{n\allowbreak{}+1\}H. Compatibility converts the second sum into one with coefficients a\_\allowbreak{}n t\_\allowbreak{}\{m\_\allowbreak{}n,\allowbreak{}r\}; subtracting the two intervals leaves the interval [m\_\allowbreak{}n,\allowbreak{}m\_\allowbreak{}\{n\allowbreak{}+1\})\allowbreak{} minus [M\_\allowbreak{}n,\allowbreak{}M\_\allowbreak{}\{n\allowbreak{}+1\})\allowbreak{},\allowbreak{} exactly a\textquotesingle\textquotesingle-b\textquotesingle\textquotesingle. Thus enlargement does not change the maps on ker Δ or coker Δ. Two representatives of the same pro-map have a common increasing enlargement where their component maps agree; (3.2)\allowbreak{} proves their maps on these groups agree.

For a second pro-map b\_\allowbreak{}n:\allowbreak{}L\_\allowbreak{}\{ℓ\_\allowbreak{}n\}\allowbreak{}→N\_\allowbreak{}n,\allowbreak{} the composite uses indices m\_\allowbreak{}\{ℓ\_\allowbreak{}n\} and components b\_\allowbreak{}n a\_\allowbreak{}\{ℓ\_\allowbreak{}n\}. The degree-zero product map is b\textquotesingle a\textquotesingle. In degree one,\allowbreak{} expand b''a'':\allowbreak{} its outer interval ℓ\_\allowbreak{}n≤q<ℓ\_\allowbreak{}\{n\allowbreak{}+1\} and inner intervals m\_\allowbreak{}q≤r<m\_\allowbreak{}\{q\allowbreak{}+1\} partition [m\_\allowbreak{}\{ℓ\_\allowbreak{}n\},\allowbreak{}m\_\allowbreak{}\{ℓ\_\allowbreak{}\{n\allowbreak{}+1\}\})\allowbreak{}. Compatibility converts each coefficient to b\_\allowbreak{}n a\_\allowbreak{}\{ℓ\_\allowbreak{}n\}t\_\allowbreak{}\{m\_\allowbreak{}\{ℓ\_\allowbreak{}n\},\allowbreak{}r\}. Hence its degree-one map is exactly the composite\textquotesingle s map. For an identity represented by transitions t\_\allowbreak{}\{n,\allowbreak{}m\_\allowbreak{}n\},\allowbreak{} compare to m\_\allowbreak{}n\allowbreak{}=n in (3.2)\allowbreak{}. The induced maps on kernels and cokernels are identities. This proves functoriality and invariance under a pro-isomorphism.

Apply this to the towers Hom(T,\allowbreak{}K\_\allowbreak{}n)\allowbreak{} and Hom(T,\allowbreak{}K\_\allowbreak{}n[-1]\allowbreak{})\allowbreak{}. A pro-isomorphism gives isomorphisms on both outer terms of (2.1)\allowbreak{}. The map of triangles supplies a commuting map of its short exact sequences,\allowbreak{} so the map Hom(T,\allowbreak{}C\_\allowbreak{}K)\allowbreak{}\allowbreak{}→Hom(T,\allowbreak{}C\_\allowbreak{}L)\allowbreak{} is an isomorphism for every T. Yoneda then makes C\_\allowbreak{}K\allowbreak{}→C\_\allowbreak{}L an isomorphism. This proves the asserted receiving More on Algebra result using the same a',\allowbreak{}a'' as the source,\allowbreak{} including any TR3 completion. It does not assert that the resulting map itself is uniquely chosen.

\subsection{4. The truncation tower and choice of its comparison}\label{proofs-10435-10686-the-truncation-tower-and-choice-of-its-comparison}

Let T\_\allowbreak{}n\allowbreak{}=τ\_\allowbreak{}\{\textbackslash{}geq -n\}K and assume its product in D(A)\allowbreak{} exists. Let C be its homotopy limit. The canonical maps c\_\allowbreak{}n:\allowbreak{}K\allowbreak{}→T\_\allowbreak{}n are compatible,\allowbreak{} so (2.1)\allowbreak{} supplies c:\allowbreak{}K\allowbreak{}→C with π\_\allowbreak{}n j c\allowbreak{}=c\_\allowbreak{}n. For each fixed p,\allowbreak{} put N\_\allowbreak{}p\allowbreak{}=max(1,\allowbreak{}-p)\allowbreak{}. The cohomology tower has H\textasciicircum{}p(T\_\allowbreak{}n)\allowbreak{}\allowbreak{}=0 for n<N\_\allowbreak{}p and H\textasciicircum{}p(T\_\allowbreak{}n)\allowbreak{}\allowbreak{}=H\textasciicircum{}p(K)\allowbreak{} for n≥N\_\allowbreak{}p,\allowbreak{} with identity transitions on the latter factors.

For m≥N\_\allowbreak{}p,\allowbreak{} the composite H\textasciicircum{}p(K)\allowbreak{}\allowbreak{}→H\textasciicircum{}p(C)\allowbreak{}\allowbreak{}→H\textasciicircum{}p(T\_\allowbreak{}m)\allowbreak{}\allowbreak{}=H\textasciicircum{}p(K)\allowbreak{} is the identity. Thus H\textasciicircum{}p(c)\allowbreak{} is an isomorphism if and only if the second map is an isomorphism:\allowbreak{} when either is invertible the other is its inverse. The second map is independent of c. Any isomorphism of choices of the triangle preserves these projections,\allowbreak{} so the criterion is also independent of the chosen homotopy limit. This proves original 10574–\allowbreak{}10585. Existing unit 0887 removes the redundant connector without changing this argument.

The comparison c is one possible lift,\allowbreak{} not a canonical choice deduced from compatible stage maps. In particular,\allowbreak{} showing that a specific chain map γ supplies one such lift proves the criterion for every lift,\allowbreak{} by the preceding argument; it does not prove that all lifts are literally equal.

\subsection{5. The split injective tower and the actual chain map}\label{proofs-10435-10686-the-split-injective-tower-and-the-actual-chain-map}

Assume countable products and enough injectives in A. Use the strictly commuting diagram of original Lemma lemma-special-inverse-system:\allowbreak{} maps τ\_\allowbreak{}\{\textbackslash{}geq -n\}K\allowbreak{}→I\_\allowbreak{}n are quasi-isomorphisms and the maps f\_\allowbreak{}\{n\allowbreak{}+1\}:\allowbreak{}I\_\allowbreak{}\{n\allowbreak{}+1\}\allowbreak{}→I\_\allowbreak{}n are termwise split epimorphisms. Each I\_\allowbreak{}n is bounded below with injective terms,\allowbreak{} hence K-injective.

Fix a degree p and choose sections s\_\allowbreak{}\{n\allowbreak{}+1\}\textasciicircum{}p of f\_\allowbreak{}\{n\allowbreak{}+1\}\textasciicircum{}p. Put C\_\allowbreak{}1\textasciicircum{}p\allowbreak{}=I\_\allowbreak{}1\textasciicircum{}p and C\_\allowbreak{}n\textasciicircum{}p\allowbreak{}=ker f\_\allowbreak{}n\textasciicircum{}p for n≥2. The successive chosen decompositions identify I\_\allowbreak{}n\textasciicircum{}p with the finite sum of C\_\allowbreak{}1\textasciicircum{}p,\allowbreak{}…,\allowbreak{}C\_\allowbreak{}n\textasciicircum{}p,\allowbreak{} keeping the transition map as the projection dropping its last component. The termwise limit is therefore I\textasciicircum{}p\allowbreak{}=∏\_\allowbreak{}n C\_\allowbreak{}n\textasciicircum{}p,\allowbreak{} with the actual projections λ\_\allowbreak{}n\textasciicircum{}p:\allowbreak{}I\textasciicircum{}p\allowbreak{}→I\_\allowbreak{}n\textasciicircum{}p from those decompositions. The differentials on the I\_\allowbreak{}n induce a unique differential on I commuting with all λ\_\allowbreak{}n; its square is zero after every projection,\allowbreak{} so is zero. Thus this is a complex and the termwise limit.

Let P\textasciicircum{}p\allowbreak{}=∏\_\allowbreak{}n I\_\allowbreak{}n\textasciicircum{}p and let Δ\^{}p have components x\_\allowbreak{}n-f\_\allowbreak{}\{n\allowbreak{}+1\}\textasciicircum{}p x\_\allowbreak{}\{n\allowbreak{}+1\}. Define a section σ\textasciicircum{}p:\allowbreak{}P\textasciicircum{}p\allowbreak{}→P\textasciicircum{}p recursively by

(σ\^{}p y)\allowbreak{}\_\allowbreak{}1\allowbreak{}=0,\allowbreak{} (σ\^{}p y)\allowbreak{}\_\allowbreak{}\{n\allowbreak{}+1\}\allowbreak{}=s\_\allowbreak{}\{n\allowbreak{}+1\}\textasciicircum{}p((σ\textasciicircum{}p y)\allowbreak{}\_\allowbreak{}n-y\_\allowbreak{}n)\allowbreak{}. (5.1)\allowbreak{}

Every component is a finite expression in the projections of y,\allowbreak{} so (5.1)\allowbreak{} defines a morphism in A by the product property. It satisfies Δ\^{}p σ\textasciicircum{}p\allowbreak{}=id exactly. The compatible sequences are exactly im λ\textasciicircum{}p,\allowbreak{} which proves

0\allowbreak{}→I\textasciicircum{}p --λ\textasciicircum{}p-\allowbreak{}-> P\^{}p --Δ\textasciicircum{}p-\allowbreak{}-> P\textasciicircum{}p\allowbreak{}→0 (5.2)\allowbreak{}

is split exact. To identify the other splitting mentioned in the source,\allowbreak{} let ρ\textasciicircum{}p:\allowbreak{}P\textasciicircum{}p\allowbreak{}→∏C\_\allowbreak{}n\textasciicircum{}p take from x\_\allowbreak{}n its C\_\allowbreak{}n\textasciicircum{}p component. Then ρ\^{}p λ\textasciicircum{}p\allowbreak{}=id,\allowbreak{} ρ\^{}p σ\textasciicircum{}p\allowbreak{}=0,\allowbreak{} and λ\textasciicircum{}pρ\textasciicircum{}p\allowbreak{}+σ\textasciicircum{}pΔ\textasciicircum{}p\allowbreak{}=id:\allowbreak{} subtract σ\^{}pΔ\^{}p x from x to get a compatible sequence,\allowbreak{} whose C\_\allowbreak{}n coordinates are precisely ρ\^{}p x.

The splittings need not be chain maps,\allowbreak{} and are not asserted to be. The λ and Δ are chain maps and give a termwise split exact sequence of complexes. Its triangle in K(A)\allowbreak{} and D(A)\allowbreak{} is I\allowbreak{}→P\allowbreak{}→P\allowbreak{}→I[1]\allowbreak{}. The product-of-K-injectives argument from PART14 identifies P as the derived product of the I\_\allowbreak{}n and makes it K-injective. Closure under triangles then makes I K-injective. Equivalently PART15\textquotesingle s split inverse-system argument proves this. The displayed triangle identifies I with holim I\_\allowbreak{}n and,\allowbreak{} through the stage quasi-isomorphisms,\allowbreak{} with holim τ\_\allowbreak{}\{\textbackslash{}geq -n\}K.

The strictly compatible composite maps K\allowbreak{}→τ\_\allowbreak{}\{\textbackslash{}geq -n\}K\allowbreak{}→I\_\allowbreak{}n give a unique actual chain map γ:\allowbreak{}K\allowbreak{}→I by the termwise limit property. It represents a lift c\_\allowbreak{}γ:\allowbreak{}K\allowbreak{}→holim τ\_\allowbreak{}\{\textbackslash{}geq -n\}K with the prescribed stage maps. Hence γ is a quasi-isomorphism exactly when c\_\allowbreak{}γ is an isomorphism. Section 4 makes this equivalent to any chosen comparison c being an isomorphism. This justifies the source\textquotesingle s conclusion without identifying all possible c with a single predetermined map in its displayed diagram. There is no new source-defect count for that choice convention.

No assertion that γ is a quasi-isomorphism has yet been made under countable products alone. Its obstruction is computed next.

\subsection{6. The exact recovery obstruction without exact products}\label{proofs-10435-10686-the-exact-recovery-obstruction-without-exact-products}

DERIVED-UNDER-025 concerns the actual truncation tower T\_\allowbreak{}n,\allowbreak{} assuming only that its product P and homotopy limit C exist in D(A)\allowbreak{}. Exactness of products in A and enough injectives are not assumed. Write A\_\allowbreak{}p\allowbreak{}=H\textasciicircum{}p(P)\allowbreak{} and δ\_\allowbreak{}p\allowbreak{}=H\textasciicircum{}p(Δ)\allowbreak{}. The triangle gives a short exact sequence in A

0\allowbreak{}→C\_\allowbreak{}p\allowbreak{}→H\textasciicircum{}p(C)\allowbreak{} --z\_\allowbreak{}p-\allowbreak{}-> Z\_\allowbreak{}p\allowbreak{}→0,\allowbreak{} C\_\allowbreak{}p\allowbreak{}=coker δ\_\allowbreak{}\{p-1\},\allowbreak{} Z\_\allowbreak{}p\allowbreak{}=ker δ\_\allowbreak{}p. (6.1)\allowbreak{}

The first injection is induced by the shifted boundary of (1.1)\allowbreak{}; thus its sign is the one specified in Section 2. All these kernels and cokernels exist since A is abelian. There is no replacement of A\_\allowbreak{}p by a product of cohomology objects.

For m≥N\_\allowbreak{}p the projection P\allowbreak{}→T\_\allowbreak{}m induces q\_\allowbreak{}p:\allowbreak{}Z\_\allowbreak{}p\allowbreak{}→H\textasciicircum{}p(K)\allowbreak{}. This map is independent of m:\allowbreak{} on Z\_\allowbreak{}p the equations δ\_\allowbreak{}p\allowbreak{}=0 say that adjacent projections are related by the cohomology transition maps,\allowbreak{} which are identities after N\_\allowbreak{}p. The canonical product map κ:\allowbreak{}K\allowbreak{}→P has components c\_\allowbreak{}n. The equality Δκ\allowbreak{}=0 lets H\textasciicircum{}p(κ)\allowbreak{} factor as s\_\allowbreak{}p:\allowbreak{}H\textasciicircum{}p(K)\allowbreak{}\allowbreak{}→Z\_\allowbreak{}p,\allowbreak{} and q\_\allowbreak{}p s\_\allowbreak{}p\allowbreak{}=id. Define B\_\allowbreak{}p\allowbreak{}=ker q\_\allowbreak{}p and

E\_\allowbreak{}p\allowbreak{}=ker(q\_\allowbreak{}p z\_\allowbreak{}p:\allowbreak{}H\textasciicircum{}p(C)\allowbreak{}\allowbreak{}→H\textasciicircum{}p(K)\allowbreak{})\allowbreak{}.

Pull back the epimorphism z\_\allowbreak{}p along B\_\allowbreak{}p\allowbreak{}→Z\_\allowbreak{}p in (6.1)\allowbreak{}. Its kernel remains C\_\allowbreak{}p and the pullback object is E\_\allowbreak{}p,\allowbreak{} so

0\allowbreak{}→C\_\allowbreak{}p\allowbreak{}→E\_\allowbreak{}p\allowbreak{}→B\_\allowbreak{}p\allowbreak{}→0 (6.2)\allowbreak{}

is exact. These are explicit objects with the indicated maps,\allowbreak{} independent of the choice of c. For any lift c with the required stage maps,\allowbreak{} (q\_\allowbreak{}p z\_\allowbreak{}p)\allowbreak{}H\textasciicircum{}p(c)\allowbreak{}\allowbreak{}=id. Consequently

H\textasciicircum{}p(K)\allowbreak{}\allowbreak{}⊕E\_\allowbreak{}p\allowbreak{}→H\textasciicircum{}p(C)\allowbreak{},\allowbreak{} (x,\allowbreak{}e)\allowbreak{}\allowbreak{}↦H\textasciicircum{}p(c)\allowbreak{}x\allowbreak{}+e (6.3)\allowbreak{}

is an isomorphism. Its inverse sends y to ((q\_\allowbreak{}p z\_\allowbreak{}p)\allowbreak{}y,\allowbreak{} y-H\textasciicircum{}p(c)\allowbreak{}(q\_\allowbreak{}p z\_\allowbreak{}p)\allowbreak{}y)\allowbreak{}. This writes out the actual splitting without suppressing the defect. It depends on the chosen c; E\_\allowbreak{}p itself does not.

It follows that every such c is an isomorphism in D(A)\allowbreak{} exactly when

coker δ\_\allowbreak{}\{p-1\}\allowbreak{}=0 and ker(q\_\allowbreak{}p:\allowbreak{}ker δ\_\allowbreak{}p\allowbreak{}→H\textasciicircum{}p(K)\allowbreak{})\allowbreak{}\allowbreak{}=0 for every integer p. (6.4)\allowbreak{}

Indeed (6.2)\allowbreak{} makes these conditions equivalent to E\_\allowbreak{}p\allowbreak{}=0,\allowbreak{} and (6.3)\allowbreak{} then makes every H\textasciicircum{}p(c)\allowbreak{} invertible. Conversely invertibility of c makes the kernel in (6.3)\allowbreak{} zero,\allowbreak{} forcing both outer objects in (6.2)\allowbreak{} to vanish. This is a proved necessary-and-sufficient criterion,\allowbreak{} not an assumption that the missing product exactness holds. Under Section 5\textquotesingle s hypotheses it is also the exact criterion for the specific γ:\allowbreak{}K\allowbreak{}→I to be a quasi-isomorphism.

\subsection{\texorpdfstring{7. Exact products,\allowbreak{} stabilization,\allowbreak{} and removal of an unused hypothesis}{7. Exact   and removal of an unused hypothesis}}\label{proofs-10435-10686-exact-products-stabilization-and-removal-of-an-unused-hypothesis}

DERIVED-UNDER-026 assumes exact countable products in A,\allowbreak{} but does not assume enough injectives. Section 1 supplies products and homotopy limits for every countable tower in D(A)\allowbreak{}. For any such tower K\_\allowbreak{}n its triangle has the cohomology exact sequence

\begingroup\raggedright
0\allowbreak{}→coker(Δ:\allowbreak{}∏H\textasciicircum{}\{p-1\}(K\_\allowbreak{}n)\allowbreak{}\allowbreak{}→∏H\textasciicircum{}\{p-1\}(K\_\allowbreak{}n)\allowbreak{})\allowbreak{} \allowbreak{}→H\textasciicircum{}p(holim K\_\allowbreak{}n)\allowbreak{} \allowbreak{}→ker(Δ:\allowbreak{}∏H\textasciicircum{}p(K\_\allowbreak{}n)\allowbreak{}\allowbreak{}→∏H\textasciicircum{}p(K\_\allowbreak{}n)\allowbreak{})\allowbreak{}\allowbreak{}→0. (7.1)\allowbreak{}
\par\endgroup

All maps are induced by the original projections,\allowbreak{} transitions and boundary of the chosen triangle. This follows from (1.1)\allowbreak{} and H\^{}p of the actual product,\allowbreak{} as proved in Section 1.

Here is the needed vanishing for any eventually isomorphic tower B\_\allowbreak{}n in A,\allowbreak{} with f\_\allowbreak{}\{n\allowbreak{}+1\} isomorphisms for n≥N. Write t\_\allowbreak{}\{r,\allowbreak{}s\}:\allowbreak{}B\_\allowbreak{}s\allowbreak{}→B\_\allowbreak{}r for its unchanged transition composites. Define σ:\allowbreak{}∏B\_\allowbreak{}n\allowbreak{}→∏B\_\allowbreak{}n by finite component formulas

(σ y)\allowbreak{}\_\allowbreak{}N\allowbreak{}=0,\allowbreak{} (σ y)\allowbreak{}\_\allowbreak{}n\allowbreak{}=Σ\_\allowbreak{}\{n≤j<N\}t\_\allowbreak{}\{n,\allowbreak{}j\}y\_\allowbreak{}j,\allowbreak{} n<N,\allowbreak{} (σ y)\allowbreak{}\_\allowbreak{}n\allowbreak{}=-Σ\_\allowbreak{}\{N≤j<n\}(t\_\allowbreak{}\{j,\allowbreak{}n\})\allowbreak{}\textasciicircum{}(-1)\allowbreak{}y\_\allowbreak{}j,\allowbreak{} n\textgreater N. (7.2)\allowbreak{}

Every inverse in the last formula is of a transition within the isomorphism range. Direct cancellation gives (σy)\allowbreak{}\_\allowbreak{}n-f\_\allowbreak{}\{n\allowbreak{}+1\}(σy)\allowbreak{}\_\allowbreak{}\{n\allowbreak{}+1\}\allowbreak{}=y\_\allowbreak{}n,\allowbreak{} also at n\allowbreak{}=N. Thus Δσ\allowbreak{}=id. Its kernel is B\_\allowbreak{}N by the map with components t\_\allowbreak{}\{n,\allowbreak{}N\} for n\textless N and (t\_\allowbreak{}\{N,\allowbreak{}n\})\allowbreak{}\textasciicircum{}(-1)\allowbreak{} for n≥N. The inverse is projection at N. These exact formulas retain the initial finite part and all transition factors.

Apply this to the truncation tower\textquotesingle s H\^{}p. Its nonzero tail starts at N\_\allowbreak{}p\allowbreak{}=max(1,\allowbreak{}-p)\allowbreak{}. The preceding calculation makes C\_\allowbreak{}p\allowbreak{}=0 in (6.1)\allowbreak{} and identifies Z\_\allowbreak{}p with H\textasciicircum{}p(K)\allowbreak{} through the very map q\_\allowbreak{}p,\allowbreak{} so B\_\allowbreak{}p\allowbreak{}=0 as well. Section 6 therefore proves every comparison

K\allowbreak{}→holim\_\allowbreak{}n τ\_\allowbreak{}\{\textbackslash{}geq -n\}K

is an isomorphism,\allowbreak{} without enough injectives. This is the asserted strengthening of the intermediate recovery claim in original 10670–\allowbreak{}10685. Enough injectives remain necessary as an assumption of the source\textquotesingle s separate construction of the I\_\allowbreak{}n; retaining them,\allowbreak{} Section 5 now gives the stated K-injective resolution. The source theorem is not replaced by a stronger existence claim.

The same proof gives a more general convergence assertion. If for a fixed p the towers H\textasciicircum{}\{p-1\}(K\_\allowbreak{}n)\allowbreak{} and H\textasciicircum{}p(K\_\allowbreak{}n)\allowbreak{} are eventually isomorphic,\allowbreak{} (7.1)\allowbreak{} identifies H\^{}p(holim K\_\allowbreak{}n)\allowbreak{} with the eventual H\textasciicircum{}p(K\_\allowbreak{}n)\allowbreak{},\allowbreak{} by the actual stable projection. In fact for this one p it suffices that the first difference map in (7.1)\allowbreak{} is epic and that the kernel of the second projects isomorphically to the chosen stable value. If E\allowbreak{}→K\_\allowbreak{}n are compatible and induce isomorphisms on H\^{}p for all sufficiently large n for each p,\allowbreak{} then any lift E\allowbreak{}→holim K\_\allowbreak{}n is an isomorphism:\allowbreak{} the compatibility makes the cohomology transitions eventually isomorphic in every degree,\allowbreak{} and the stable projection composite is exactly the given isomorphism. The stabilization index may depend on p; no uniform bound is used.

Existing unit MC-STK-ERR-0882 makes n the explicit product index in original 10681 while keeping the condition p≥-n. For fixed p it is precisely the product over n≥N\_\allowbreak{}p used here. No bound on the fixed degree p is substituted for the tower index.

\subsection{8. The source\textquotesingle s omitted dual statements}\label{proofs-10435-10686-the-sources-omitted-dual-statements}

The source says that the duals of its three colimit lemmas could be stated here. They follow with the following exact hypotheses. These supplementary proofs are not counted as source errors.

First,\allowbreak{} assume all countable inverse limits in A exist and the inverse-limit functor on towers is exact. The inverse limit of the finite products ∏\_\allowbreak{}\{j≤n\}A\_\allowbreak{}j,\allowbreak{} with the projection transitions,\allowbreak{} is the countable product by its universal property. Finite products are exact,\allowbreak{} so exact inverse limits make countable products exact. For a tower A\_\allowbreak{}n,\allowbreak{} consider the finite split exact sequence

0\allowbreak{}→A\_\allowbreak{}n\allowbreak{}→∏\_\allowbreak{}\{j≤n\}A\_\allowbreak{}j\allowbreak{}→∏\_\allowbreak{}\{j<n\}A\_\allowbreak{}j\allowbreak{}→0.

The first map has components t\_\allowbreak{}\{j,\allowbreak{}n\}; the second has components x\_\allowbreak{}j-f\_\allowbreak{}\{j\allowbreak{}+1\}x\_\allowbreak{}\{j\allowbreak{}+1\} for j\textless n. Its kernel is the image of the first map,\allowbreak{} and it is split epic:\allowbreak{} prescribe x\_\allowbreak{}n\allowbreak{}=0 and recursively solve x\_\allowbreak{}j\allowbreak{}=y\_\allowbreak{}j\allowbreak{}+f\_\allowbreak{}\{j\allowbreak{}+1\}x\_\allowbreak{}\{j\allowbreak{}+1\} backwards. The sequence maps to the preceding one by the projections on the products and the transition A\_\allowbreak{}n\allowbreak{}→A\_\allowbreak{}\{n-1\}. Its exact inverse limit is

0\allowbreak{}→lim A\_\allowbreak{}n\allowbreak{}→∏A\_\allowbreak{}n --Δ-\allowbreak{}-> ∏A\_\allowbreak{}n\allowbreak{}→0. (8.1)\allowbreak{}

For a tower of complexes,\allowbreak{} apply (8.1)\allowbreak{} termwise. The short exact sequence yields a distinguished triangle in D(A)\allowbreak{},\allowbreak{} and Section 1 identifies its two products as derived products. Thus the termwise inverse limit is a homotopy limit. Exact countable products alone are not substituted for exactness of the inverse-limit functor in this statement.

Finally let D have countable products and let Q satisfy \allowbreak{}⊕Hom\_\allowbreak{}D(E\_\allowbreak{}n,\allowbreak{}Q)\allowbreak{}≅Hom\_\allowbreak{}D(∏E\_\allowbreak{}n,\allowbreak{}Q)\allowbreak{} for every countable family. For any inverse tower with limit triangle (1.1)\allowbreak{},\allowbreak{} applying Hom\_\allowbreak{}D(-,\allowbreak{}Q)\allowbreak{} gives

Hom(P,\allowbreak{}Q)\allowbreak{} --Δ\textasciicircum{}*-\allowbreak{}-> Hom(P,\allowbreak{}Q)\allowbreak{}\allowbreak{}→Hom(C,\allowbreak{}Q)\allowbreak{}\allowbreak{}→ Hom(P[-1]\allowbreak{},\allowbreak{}Q)\allowbreak{} --Δ[-1]\allowbreak{}\textasciicircum{}*-\allowbreak{}-> Hom(P[-1]\allowbreak{},\allowbreak{}Q)\allowbreak{}.

Under the stipulated direct-sum identifications,\allowbreak{} both difference maps have the finite-support columns 1-f\_\allowbreak{}\{n\allowbreak{}+1\}\textasciicircum{}* and are injective. Their cokernel in the first row is colim Hom(K\_\allowbreak{}n,\allowbreak{}Q)\allowbreak{}. Hence the canonical map colim Hom(K\_\allowbreak{}n,\allowbreak{}Q)\allowbreak{}\allowbreak{}→Hom(C,\allowbreak{}Q)\allowbreak{} is an isomorphism. The injection proof is by successive finite-support coordinates,\allowbreak{} exactly as in PART16. This proves the final dual statement without invoking any hypothetical exact inverse-limit functor on D.

\subsection{9. Receiving checks and disposition totals}\label{proofs-10435-10686-receiving-checks-and-disposition-totals}

More on Algebra current 23752–\allowbreak{}23824 and 23828–\allowbreak{}23846 receive the explicit pro-map homotopies and Hom exact sequence in Sections 2–\allowbreak{}3. Its later pro-isomorphism consequence is thereby proved with the actual maps a',\allowbreak{}a'' and compatible boundaries. The repeated spelling at current 23772 and the abbreviated transition notation at current 23820 are routed for that chapter\textquotesingle s own report comparison,\allowbreak{} not newly adjudicated here.

More on Algebra current 14522–\allowbreak{}14533 has both exact products and enough injectives for modules,\allowbreak{} so the K-injective existence and derived-limit invocations follow exactly from Sections 1,\allowbreak{}5,\allowbreak{}7. The ring-module hypotheses were not weakened in this receiving check.

Cohomology current 10338–\allowbreak{}10341 and Cohomology on Sites current 5604–\allowbreak{}5607 use the equivalence between the particular injective tower\textquotesingle s quasi-isomorphism and truncation recovery,\allowbreak{} as proved in Sections 4–\allowbreak{}5. Perfect Complexes current 138–\allowbreak{}143 uses that equivalence in its footnote. These uses are valid; their separate geometric recovery hypotheses are not replaced by an assertion that products of arbitrary sheaves are exact. Cohomology current 10087–\allowbreak{}10116 explicitly verifies local recovery through stable cohomology maps. Section 6 describes the exact global defect those hypotheses must kill,\allowbreak{} while preserving that chapter\textquotesingle s geometric proof for its own review. The product side of Duality current 798–\allowbreak{}801,\allowbreak{} queued in PART16,\allowbreak{} is now justified by Section 1:\allowbreak{} for the shifts M\_\allowbreak{}n[n]\allowbreak{} there is only one nonzero term in each degree,\allowbreak{} and modules have exact products.

Five physical reports are complete in new groups DERIVED-RECON-081,\allowbreak{} 082 and 083:\allowbreak{} two reports each for units 0886 and 0887,\allowbreak{} and one for 0882. Totals are 142/\allowbreak{}176 reports,\allowbreak{} 83 groups,\allowbreak{} 131 already corrected,\allowbreak{} ten missing copyedit/\allowbreak{}notation reports and one missing mathematical-index report. There are 76 prior Derived units and 93 prior operations reconciled. This pass adds no new correction operation or unreported defect. The cumulative 69 proposed operations and 36 unreported findings remain unchanged. UNDER-025 and UNDER-026 bring the editorial consequence count to 26.

Review is complete through original 10686. Reading has reached 10840; the following operations-on-subcategories section is read ahead only. Next semantic review starts at 10691; unread source starts at 10841. The correction queue remains first,\allowbreak{} and no new session,\allowbreak{} live source composition,\allowbreak{} TeX or Lean run,\allowbreak{} cleanup,\allowbreak{} publication or producer mathematical acceptance occurred.
